\Askhsh{31746}{4}
\begin{ylh}{1}
    \item 2.2 Παραγωγίσιμες συναρτήσεις - Παράγωγος συνάρτηση
    \item 2.3 Κανόνες παραγώγισης
    \item 2.6 Συνέπειες του Θεωρήματος Μέσης Τιμής
    \item 2.7 Τοπικά ακρότατα συνάρτησης
    \item 2.8 Κυρτότητα - Σημεία καμπής συνάρτησης
    \item 2.9 Ασύμπτωτες - Κανόνες De l’ Hospital.
\end{ylh}
Δίνεται η συνάρτηση $f(x)=(x^2-4x+6)e^x, x \in \mathbb{R}$.
\begin{alist}
    \item Να μελετήσετε τη συνάρτηση $f$ ως προς τη μονοτονία και να βρείτε το σύνολο τιμών της. \monades{9}
    \item Να βρείτε την εφαπτομένη της γραφικής παράστασης της συνάρτησης $f$ στο σημείο της $M(0, f(0))$. \monades{5}
    \item Να μελετήσετε τη συνάρτηση $f$ ως προς την κυρτότητα και τα σημεία καμπής. \monades{6}
    \item Να αποδείξετε ότι: $f(x) \ge 2x+6$ για κάθε $x \in \mathbb{R}$. \monades{5}
\end{alist}